\pandoclink{detailed-diagnostics-for-binding-and-cpp-baselines}{%
\section{Detailed diagnostics for binding and CPP baselines}\label{detailed-diagnostics-for-binding-and-cpp-baselines}}

The following are transcriptions of committed result tables, not new model runs. A selected 6,000-row head-to-head panel and the 16,870-row peptide-held-out benchmark are different datasets; the 12-allele G3 transfer challenge is different again. Each denominator stays attached to its own metric.

\pandoclink{within-panel-comparison-by-allele}{%
\subsection{Within-panel comparison by allele}\label{within-panel-comparison-by-allele}}

\texttt{results/per\_allele\_headtohead.json} contains 30 allele-specific comparisons on the selected same-row head-to-head panel. Overall 21 of 30 favor the ensemble numerically, but both allele prevalence and row counts vary. A per-allele numerical gain does not ensure a held-out-allele gain, as the G3 failure demonstrates.

\begin{longtable}[]{@{}lrrrrr@{}}
\toprule
Allele & Rows & Binders & MHCflurry AUC & Ensemble AUC & Difference\tabularnewline
\midrule
\endhead
HLA-A*01:01 & 136 & 102 & 0.9478 & 0.9230 & -0.0248\tabularnewline
HLA-A*02:01 & 736 & 660 & 0.9418 & 0.9420 & +0.0002\tabularnewline
HLA-A*02:02 & 229 & 218 & 0.9133 & 0.9333 & +0.0200\tabularnewline
HLA-A*02:03 & 302 & 276 & 0.9242 & 0.9331 & +0.0089\tabularnewline
HLA-A*02:06 & 262 & 241 & 0.8718 & 0.9176 & +0.0458\tabularnewline
HLA-A*03:01 & 372 & 314 & 0.8827 & 0.8731 & -0.0096\tabularnewline
HLA-A*11:01 & 341 & 296 & 0.9273 & 0.9101 & -0.0172\tabularnewline
HLA-A*23:01 & 74 & 63 & 0.9365 & 0.9380 & +0.0014\tabularnewline
HLA-A*24:02 & 128 & 104 & 0.8678 & 0.8702 & +0.0024\tabularnewline
HLA-A*24:03 & 54 & 45 & 1.0000 & 0.9753 & -0.0247\tabularnewline
HLA-A*26:01 & 89 & 59 & 0.8712 & 0.8842 & +0.0130\tabularnewline
HLA-A*29:02 & 96 & 81 & 0.9012 & 0.9267 & +0.0255\tabularnewline
HLA-A*30:01 & 99 & 88 & 0.8740 & 0.9349 & +0.0610\tabularnewline
HLA-A*31:01 & 279 & 236 & 0.9031 & 0.9225 & +0.0194\tabularnewline
HLA-A*33:01 & 142 & 114 & 0.9583 & 0.9558 & -0.0025\tabularnewline
HLA-A*68:01 & 254 & 229 & 0.8376 & 0.9027 & +0.0652\tabularnewline
HLA-A*68:02 & 226 & 186 & 0.9148 & 0.9113 & -0.0035\tabularnewline
HLA-A*69:01 & 63 & 39 & 0.9818 & 0.9562 & -0.0256\tabularnewline
HLA-B*07:02 & 196 & 171 & 0.9415 & 0.9331 & -0.0084\tabularnewline
HLA-B*08:01 & 126 & 103 & 0.8611 & 0.9164 & +0.0553\tabularnewline
HLA-B*15:01 & 195 & 166 & 0.8880 & 0.9094 & +0.0214\tabularnewline
HLA-B*15:17 & 69 & 55 & 0.9818 & 0.9909 & +0.0091\tabularnewline
HLA-B*18:01 & 53 & 27 & 0.9145 & 0.9359 & +0.0214\tabularnewline
HLA-B*27:05 & 85 & 58 & 0.9055 & 0.9623 & +0.0568\tabularnewline
HLA-B*39:01 & 44 & 33 & 0.8650 & 0.8788 & +0.0138\tabularnewline
HLA-B*40:01 & 92 & 74 & 0.9872 & 0.9910 & +0.0038\tabularnewline
HLA-B*44:02 & 68 & 56 & 0.9420 & 0.9896 & +0.0476\tabularnewline
HLA-B*51:01 & 70 & 47 & 0.8483 & 0.8085 & -0.0398\tabularnewline
HLA-B*57:01 & 76 & 58 & 0.9665 & 0.9933 & +0.0268\tabularnewline
HLA-B*58:01 & 122 & 99 & 0.9095 & 0.9368 & +0.0272\tabularnewline
\bottomrule
\end{longtable}

Comparisons on small or binder-heavy alleles can have few negatives, so the sampling variability of AUROC is uneven. The Wilcoxon result in the companion analysis summarizes the signs but is not a cluster-robust inference over independent cohorts. {[}PENDING: replicate an allele-held-out benchmark without selecting alleles or thresholds after these results.{]}

\pandoclink{peptide-length-and-assay-composition-sensitivity}{%
\subsection{Peptide-length and assay-composition sensitivity}\label{peptide-length-and-assay-composition-sensitivity}}

\texttt{results/length\_analysis.json} records the underlying peptide-length composition, heavily concentrated at length nine. \texttt{results/per\_length\_headtohead.json} narrows to the selected 6,000-row comparison and reports AUROC differences with bootstrap intervals. The two denominators must not be mixed.

\begin{longtable}[]{@{}rrrrrrl@{}}
\toprule
Length & Head-to-head rows & Binders & Ensemble AUC & MHCflurry AUC & Delta & 95\% interval\tabularnewline
\midrule
\endhead
10 & 1600 & 1410 & 0.8928 & 0.8703 & +0.0224 & {[}+0.0027, +0.0434{]}\tabularnewline
11 & 36 & 30 & 0.8500 & 0.9056 & -0.0580 & {[}-0.2857, +0.1010{]}\tabularnewline
8 & 28 & 23 & 0.8783 & 0.5826 & +0.2977 & {[}+0.0267, +0.5925{]}\tabularnewline
9 & 4330 & 3588 & 0.9392 & 0.9345 & +0.0046 & {[}-0.0029, +0.0120{]}\tabularnewline
\bottomrule
\end{longtable}

The 8-mer stratum has only 28 rows and 11-mer has 36; no mechanistic length preference should be asserted from either. Length nine has 4,330 rows and a delta interval crossing zero. Length ten has 1,600 rows and a positive within-panel delta, but it is a post-hoc stratum and does not reverse the G3 transfer failure.

\pandoclink{cpp-learning-curve-and-compositional-feature-audit}{%
\subsection{CPP learning curve and compositional feature audit}\label{cpp-learning-curve-and-compositional-feature-audit}}

The committed \texttt{results/cpp\_learning\_curve.json} fits a three-mer logistic model at several training fractions and scores the same 183-row fixed test. Four recorded seeds are used below; at the full fraction they produce identical AUROC. This is a data-efficiency diagnostic for this one split, not independent test-set replication.

\begin{longtable}[]{@{}rrrr@{}}
\toprule
Training fraction & Training sequences & Mean AUROC & Across-seed SD\tabularnewline
\midrule
\endhead
0.10 & 92 & 0.8092 & 0.0289\tabularnewline
0.25 & 230 & 0.8655 & 0.0070\tabularnewline
0.50 & 461 & 0.9030 & 0.0059\tabularnewline
0.75 & 692 & 0.9177 & 0.0040\tabularnewline
1.00 & 922 & 0.9265 & 0.0000\tabularnewline
\bottomrule
\end{longtable}

The 0.10-to-1.00 training-size increase moves mean AUROC from 0.8092 to 0.9265; it cannot by itself prove that the CNN would catch up with enough data, because the learning curve here belongs to the logistic baseline only. \texttt{results/cpp\_kmer\_feature\_weights.json} confirms exact refit-test AUROC 0.926523 and records 8,000 vocabulary features. Positive weights are strongest for short basic motifs; these are correlational model coefficients, not experimental membrane-transport mechanisms.

\begin{longtable}[]{@{}lr@{}}
\toprule
Positive 3-mer & Logistic coefficient\tabularnewline
\midrule
\endhead
\texttt{RRR} & +3.568\tabularnewline
\texttt{KKK} & +1.650\tabularnewline
\texttt{KKR} & +1.422\tabularnewline
\texttt{RKK} & +1.267\tabularnewline
\texttt{KRR} & +0.970\tabularnewline
\texttt{KLA} & +0.914\tabularnewline
\texttt{QRR} & +0.857\tabularnewline
\texttt{ARR} & +0.833\tabularnewline
\texttt{KRK} & +0.824\tabularnewline
\texttt{LAL} & +0.821\tabularnewline
\texttt{RRQ} & +0.794\tabularnewline
\texttt{RQR} & +0.789\tabularnewline
\texttt{LAK} & +0.765\tabularnewline
\texttt{LLK} & +0.726\tabularnewline
\texttt{RRA} & +0.681\tabularnewline
\bottomrule
\end{longtable}

Of the top 50 positive terms, 46 contain R or K, four contain W/Y/F, and 18 contain L/I/V according to the committed composition summary. Feature frequency and coefficient signs depend on the training set and background negatives; they may reflect source and amino-acid composition rather than true uptake. The separate CPP screen still needs measured intracellular delivery and hemolysis.

\pandoclink{interpretation}{%
\subsection{Interpretation}\label{interpretation}}

These detailed tables explain two apparent contradictions without resolving them by rhetoric: the ensemble can win narrowly on a selected familiar-allele panel while failing badly on novel alleles; and a transparent k-mer baseline can beat a CNN on a small filtered CPP set while still being a poor proxy for measured cell entry. The next contribution must come from a new locked evaluation or wet-lab data, not extra pages of favorable post-hoc slices.
